% !TeX root = ../main.tex

\chapter{总结与展望}\label{chap:7}

第1章与第2章分别完成研究问题界定与文献背景梳理，第3章至第6章围绕“在保持线性推理复杂度 $O(n\bar{k}d)$ 的前提下，以可计算图结构特征驱动节点自适应稀疏分解”这一主线，依次完成问题形式化与指标提取（第3章）、三维自适应参数机制设计（第4章）、算法整合与误差分析（第5章）及实验评估（第6章）。本章在此基础上进行整体评价：第~\ref{sec:ch7-summary}~节回顾各章核心工作；第~\ref{sec:ch7-contributions}~节归纳主要创新贡献；第~\ref{sec:ch7-limitations}~节评估适用边界；第~\ref{sec:ch7-future}~节给出对应的后续方向。


\begin{figure}[htbp]
  \centering
            \includegraphics[width=1\linewidth]{7-1.pdf}
  % \setlength{\fboxsep}{1.2em}
  % \fbox{\parbox{0.88\textwidth}{\small\centering
  %   \textbf{【图7.1 占位符，定稿前须替换为正式图】}\\[0.6em]
  %   图题：MSAS-GNN 谱驱动自适应稀疏化框架贡献结构图。\\[0.4em]
  %   表达内容：以纵向层次结构展示从问题输入到实验验证的研究链条。底层标注推理复杂度约束 $O(n\bar{k}d)$；第二层展示第3章四类图复杂度指标及其代表量（$E_{\text{spectral}}(i)$、$H_{\text{norm}}(i)$、$C_{\text{deg}}(i)$、$\text{core}(i)$、$\lambda_{\text{gap}}$）；第三层展示第4章三维自适应参数（$w_k(i)$、$\tau(i)$、$k_i^{(l)}$）及第4.1--4.3节的理论分析链条；第四层展示第5章算法5.1/5.2与谱滤波误差界（式~\eqref{eq:ch5-tau-grace}）、条件化联合界（式~\eqref{eq:ch5-budget-lagrange}）；顶层标注第6章实验验证（主实验、消融、效率）。
  % }}
  \caption{MSAS-GNN 谱驱动自适应稀疏化框架贡献结构}
  \label{fig:ch7-framework}
\end{figure}


\section{研究工作总结}\label{sec:ch7-summary}

\subsection{问题背景与核心挑战}\label{subsec:ch7-background}

本文所针对的核心问题不是再次论证图神经网络在线推理为何重要，而是回答：在保持 SDGNN 线性推理主项复杂度 $O(n\bar{k}d)$ 的前提下，如何利用可计算图结构特征实现节点级差异化稀疏分配，并在仅引入可接受常数开销的情况下提升准确性与结构适应能力。围绕这一问题，后续各章依次完成了指标提取、参数设计、算法整合与实验检验。

\subsection{各章核心工作回顾}\label{subsec:ch7-chapters}

第1章完成研究问题界定与研究边界说明，第2章建立谱图理论、高效稀疏 GNN 与复杂图结构学习的背景基础；在此基础上，第3章建立问题形式化框架，经 Lanczos 谱计算提取四类图复杂度指标——谱间隙 $\lambda_{\text{gap}}$、节点谱能量 $E_{\text{spectral}}(i)$、归一化局部图熵 $H_{\text{norm}}(i)$ 与节点中心性（$C_{\text{deg}}(i)$、$\text{core}(i)$）——并将其通过明确接口传递至第4章三个参数维度。

第4章将上述指标转化为三维自适应参数。第4章第4.1节在保持任务通用性的前提下采用等权实例化构造节点谱能量 $E_{\text{spectral}}(i)$；节点维度依据谱扰动分析的基本结论，构造以节点谱能量为核心输入的节点正则系数 $\tau(i)$，建立“谱能量大$\Rightarrow\tau(i)$小$\Rightarrow$保留邻居多”的单调映射；第4章第4.3节依据多跳传播信息逐层衰减的经验规律，设计分层预算分配策略，给出 $k_i^{(l)}$ 的分配参考形式。如此构成三维参数体系，输出至第5章进行端到端整合。

第5章完成算法整合与误差分析：算法~5.1 以 $k_i^{(l)}$ 为逐层预算约束、$\tau(i)$ 为正则系数驱动 LARS 分层求解；算法~5.2 固化 $\Theta^{\text{fixed}}$，推理主项复杂度维持 $O(n\bar{k}d)$。第5章第5.2节分别给出多项式谱滤波近似误差上界（式~\eqref{eq:ch5-tau-grace}）与条件化误差界（式~\eqref{eq:ch5-budget-lagrange}），为第6章谱相似性验证实验提供分析接口。

第6章以引文网络、网页图与大规模图上的主实验、逐模块消融、参数敏感性分析和效率实测验证前三章主张。网页图上的更高精度增益说明拓扑自适应稀疏预算在结构更复杂的数据上收益更明显；在与 SDGNN 相同推理复杂度口径下，三维参数体系带来精度提升，说明效率与准确性可在统一框架内共同优化。

\subsection{方法、分析与实验结果的一致性说明}\label{subsec:ch7-loop}

第4章第4.2节的 $\tau(i)$ 设计准则直接进入算法~5.1 的节点级正则项，第6章消融结果支持了谱能量在参数分配中的主导作用；第4章第4.3节的衰减界分析与预算分配策略为跳距预算分配提供了方法动机与解析依据，第6章相应消融结果与该方向一致。对于第5章的谱滤波误差界（式~\eqref{eq:ch5-tau-grace}），第6章当前主要给出 $\sigma$-谱相似性随稀疏预算变化的方向性证据，而不作严格的定量闭合性验证；条件化联合误差界（式~\eqref{eq:ch5-budget-lagrange}）则主要作为复杂度与误差分析的参考支撑。上述事实共同界定了本文当前“方法设计—分析说明—实验结果”之间的一致性范围，也自然引出后文的局限性与展望。


\section{主要贡献与创新点}\label{sec:ch7-contributions}

本文的贡献可从方法设计、算法整合与实验验证三个方面加以归纳。在方法与分析贡献方面，本文从可计算图复杂度指标出发，提出了一套基于谱指标的自适应稀疏参数构造框架：频率维度在保持任务通用性的前提下采用等权实例化构造节点谱能量 $E_{\text{spectral}}(i)$，节点维度依据谱扰动分析的基本结论设计节点级正则系数 $\tau(i)$，跳距维度依据分层信息衰减规律设计分层预算分配策略；第5章第5.2节进一步给出多项式谱滤波近似误差界，并说明参数设计与稀疏误差、复杂度之间的关系。需要说明的是，本文各维度设计建立在谱分析的直觉基础上，其严格理论最优性有待后续研究进一步验证；$\tau(i)$ 设计准则更适合作为工程性参数构造理解。

算法~5.1/5.2 在统一框架内完成三维参数整合、LARS 端到端求解与推理期固化，推理主项复杂度保持 $O(n\bar{k}d)$，使精度与效率在统一口径下得以协同优化。第6章实验按精度、模块贡献与效率三类组织，统一复杂度口径下的结果表明，三维参数体系在保持同一复杂度阶的前提下以有限常数开销换取了稳定的精度收益，论证链具有清晰的可追溯性。


\section{研究局限性}\label{sec:ch7-limitations}

理论层面，第4章第4.2节的节点正则系数设计建立在谱扰动分析直觉的基础上，包含谱能量代理替换、扰动局部化等工程近似；当图结构更复杂或度分布呈强幂律时，这些近似的适用性有待进一步验证。其余近似还涉及固定度矩阵近似、Lasso 凸性条件以及从局部界到全局谱相似性的聚合近似。第5章第5.2节的谱误差界（式~\eqref{eq:ch5-tau-grace}、式~\eqref{eq:ch5-budget-lagrange}）仅在三维参数满足特定联合约束时成立，该条件在极端图上的可验证性有限。

方法层面，关键超参数（$\tau_{\text{base}}$、$E_{\text{threshold}}$）依赖网格搜索，在大规模图上代价较高，自动化调优机制尚未纳入框架。$\Theta^{\text{fixed}}$ 的固化机制使 inductive 能力受限，对新增节点需局部重优化。现有理论推导主要基于单关系静态图设定；对于异构图、多关系图及更复杂时空图场景，其适用性仍有待系统扩展。

实验层面，对于式~\eqref{eq:ch5-tau-grace}所对应的 $\sigma$-谱相似性结论，第6章当前仅给出方向性证据，尚未形成严格统一口径下的定量验证。除此之外，实验仍主要集中于引文网络与维基页面协作网络，生物信息学图、金融风控图等领域缺乏验证；可扩展性目前仅覆盖至 ogbn-arxiv 量级；部分近年高效图学习基线，尤其是图 Transformer 一类方法（如 Exphormer\cite{Shirzad2023Exphormer}、NodeFormer\cite{Wu2022NodeFormer}），尚未系统纳入比较。


\section{未来研究方向}\label{sec:ch7-future}

后续研究可进一步推进第4章第4.2节 $\tau(i)$ 设计准则的理论严格化，例如引入逐节点 Davis-Kahan 型局部化分析替代当前的全局扰动估计，使该设计准则从工程性构造演进为带有可核验条件的更严格分析结论；辅以集中不等式对幂律图上特征向量分量分布的概率刻画，可进一步覆盖极端异配与重尾度分布场景。

在方法拓展方向，超参数自动化方面可引入基于代理模型的贝叶斯优化以替代网格搜索，将 $\tau_{\text{base}}$ 纳入端到端可微优化（双层优化架构）有望使其与模型参数协同学习。异构图扩展需为不同关系类型分别构造类型感知谱能量并设计独立稀疏预算映射，是推广至知识图谱与生物医学异构网络的必要前提。inductive 场景的增强则需探索测试阶段快速适配新增节点稀疏权重的局部优化策略。

应用与系统优化方面，蛋白质相互作用网络因兼具噪声相互作用与复杂拓扑，可作为验证谱稀疏化适用边界的优先场景。量化感知训练与 LARS 定点运算优化有望推进框架向边缘部署落地；后续评测还应同步纳入更充分的最新高效 GNN 基线并统一超参数搜索预算，进一步增强结论的外部可信度。

% \vspace{0.5em}

总体来看，本文以谱分析为工具，以可计算图复杂度指标为桥梁，构建了节点自适应稀疏 GNN 推理的方法框架。现有结果在单关系静态图场景下具有较为明确的理论依据与实验支撑，并在引文网络、网页图与大规模图等不同静态场景上给出了系统实验结果；自适应稀疏化的核心思路——以拓扑复杂度指标驱动差异化资源分配——具有推广至其他图算法的潜力，而上述局限性的明确标注与后续方向的具体指向，也为该研究线索的持续深化提供了清晰起点。
